\section*{الفصل \ref{ch:g9:gravitation} --- الجاذبية}

\begin{solution}{exo:g9:gravitation:1}
$F = G \, m_A m_B / d^2$، والكتلتان \omterm{def:g3:measuring-mass:units}{بالكيلوغرام}، و$d$ \omterm{def:g2:measuring-length:units}{بالمتر}
بين المركزين، و$F$ بالنيوتن، و$G = \qty{6.67e-11}{N.m^2/kg^2}$.
وبالكلام: الجذب يكبر مع كل \omterm{def:g3:measuring-mass:mass}{كتلة} ويموت مع مربّع المسافة ---
وهو يؤثّر بالتساوي في الشريكين، في جهتين متعاكستين.
\end{solution}

\begin{solution}{exo:g9:gravitation:2}
\omterm{def:g2:pushes-and-pulls:force}{بقوة} \qty{1}{N} بالضبط --- فجذوب القانون تجيء دائمًا أزواجًا
متساوية متعاكسة. والتفاحة تتسارع تسارعًا مرئيًّا لأن كتلتها
ضئيلة؛ أما \qty{1}{N} نفسها مطبَّقة على \qty{5.97e24}{kg} من
\omterm{def:g4:solar-system:planet}{الكوكب} فتحرّكه قدرًا لا يُقاس.
\end{solution}

\begin{solution}{exo:g9:gravitation:3}
عند $2d$: ربع \omterm{def:g2:pushes-and-pulls:force}{القوة} عند $d$ --- مربّع الاثنين. وعند $10d$: جزء
من مئة.
\end{solution}

\begin{solution}{exo:g9:gravitation:4}
$F = \num{6.67e-11} \times 80 \times 80 \div 0.25 =
\qty{1.7e-6}{N}$ --- أي دون جزأين من مليون من النيوتن: فمهما
يكن الذي يقرّب شريكَي الرقص، فليس هو \omterm{def:g9:gravitation:gravitation}{الجاذبية}.
\end{solution}

\begin{solution}{exo:g9:gravitation:5}
المعدّل الذي تسقط به القذيفة، والمعدّل الذي ينحني به سطح الأرض
المستدير هابطًا تحتها. \omterm{def:g6:sun-earth-moon:axisorbit}{والمدار}: سقوط بلا وصول.
\end{solution}

\begin{solution}{exo:g9:gravitation:6}
عند ارتفاع المحطّة تحفظ \omterm{def:g9:gravitation:gravitation}{الجاذبية} قوّتها السطحية كلها تقريبًا
--- وهي التي تمسك المحطّة على دائرتها. والمحطّة وطاقمها كلاهما
\emph{يسقط}، معًا، حول الأرض: والرفقاء الساقطون لا يضغطون على
شيء، \omterm{def:g1:floating-and-sinking:words}{والطفو} هو إحساس السقوط الحرّ المشترك.
\end{solution}

\begin{solution}{exo:g9:gravitation:7}
لأن السماوات تتعامل بكتل وحشية --- فالكواكب والنجوم تضرب $G$
الواهنة فتصير قوى جبّارة --- ولأن \omterm{def:g9:gravitation:gravitation}{الجاذبية} لا تُلغى أبدًا:
فهي لا تفعل إلا الجذب، فكل \omterm{def:g3:measuring-mass:mass}{كتلة} مضافة تعمّق الشدّ.
\end{solution}

\begin{solution}{exo:g9:gravitation:8}
جذب \omterm{def:g2:moon-in-the-sky:moon}{القمر} يكوّم البحار نحوه (فالماء الأقرب مجذوب أشدّ) وبعيدًا
عنه (فالماء الأبعد مجذوب أقلّ)؛ والأرض الدائرة تحمل كل ساحل
تحت الكومتين يوميًّا --- فمدّان. وأعظمها عند المحاق \omterm{ex:g2:moon-in-the-sky:shapes}{والبدر}،
حين تصطفّ يد الشمس مع يد \omterm{def:g2:moon-in-the-sky:moon}{القمر}: مدّ الربيع.
\end{solution}

\begin{solution}{exo:g9:gravitation:9}
$F = \num{6.67e-11} \times (\num{7.3e22} \times
\num{5.97e24}) \div (\num{3.8e8})^2 \approx \qty{2.0e20}{N}$
--- أي $2$ يتبعه عشرون صفرًا من النيوتنات تمسك \omterm{def:g2:moon-in-the-sky:moon}{القمر} على
عقده.
\end{solution}

\begin{solution}{exo:g9:gravitation:10}
قرصان يديران أحدهما ضدّ الآخر: \omterm{def:g3:measuring-mass:mass}{فكتلة} \omterm{def:g2:moon-in-the-sky:moon}{القمر}، الأصغر بنحو ثمانين
مرة، تقصّ الجذب ثمانين ضعفًا؛ لكن نصف قطره، الأصغر بنحو أربع
مرات، \emph{يرفع} الجذب السطحي بمربّع الأربعة --- أي نحو أربعة
عشر ضعفًا. ثمانون هبوطًا وأربعة عشر صعودًا: فيبقى ما يقارب
السدس.
\end{solution}

\begin{solution}{exo:g9:gravitation:11}
مربّع المسافة يرقّق قبضة الشمس مع البعد، والقبضة الأضعف لا
تمسك إلا سقوطًا دائريًّا أبطأ: فالكواكب الخارجية مربوطة برخاوة
فتتمهّل، والداخلية ممسوكة بشدّة فتركض --- وعَدْو عطارد ومشية
نبتون البالغة $164$ سنة هما القانون نفسه مقروءًا عند
مسافتين.
\end{solution}

\begin{solution}{exo:g9:gravitation:12}
يجب أن يكون دوره يومًا واحدًا بالضبط، دائرًا نحو الشرق --- في
جهة دوران الأرض نفسها. فيدور عندئذٍ \omterm{def:g2:moon-in-the-sky:moon}{القمر} الصناعي والهوائي على
خطوة تامّة: فالطبق ``المتعلّق'' يراقب قذيفة مدفع تسقط حول
\omterm{def:g4:solar-system:planet}{الكوكب} بالمعدّل نفسه الذي يدور به \omterm{def:g4:solar-system:planet}{الكوكب} تحتها --- ساكنة في
عيون الراصد الدائر لا غير.
\end{solution}

\begin{solution}{pb:g9:gravitation:1}
\textbf{1.} $F = \num{6.67e-11} \times (0.1 \times
\num{5.97e24}) \div (\num{6.4e6})^2 \approx \qty{0.97}{N}$
--- ولنقل نيوتنًا واحدًا: شدّ البستان، محسوبًا من ثابت
الكون.
\textbf{2.} \qty{0.97}{N} كذلك --- التبادل: فجذوب القانون لا
تجيء إلا أزواجًا متساوية متعاكسة.
\textbf{3.} مضاعفة مسافة المركزين تربّع الجذب إلى الربع: نحو
\qty{0.24}{N}.
\textbf{4.} كل فتات من \omterm{def:g4:solar-system:planet}{الكوكب} يجذب التفاحة --- الجبال
القريبة، واللبّ العميق، والنقاط المقابلة البعيدة --- ومجموع
تلك الجذوب كلها يخرج تمامًا كأن \omterm{def:g3:measuring-mass:mass}{الكتلة} كلها جالسة في المركز:
مبرهنة لا افتراض، وقد أخّر صاحبها النشر سنوات حتى استطاع
إثباتها.
\textbf{5.} نحو $2000 \times 9.8 \approx \qty{2.0e4}{N}$ ---
أي عشرون ألف نيوتن (والقانون مع $d$ يساوي نصف قطر أرضي يعطي
النتيجة نفسها).
\textbf{6.} عند نصفَي قطر، يترك مربّع المسافة ربعًا: نحو
\qty{4900}{N} من وزنه السطحي البالغ \qty{19600}{N}.
\textbf{7.} نعم --- فربع \omterm{def:g9:gravitation:gravitation}{الجاذبية} الكاملة ليس خواءً بلا \omterm{def:g4:weight-and-mass:weight}{وزن}.
\omterm{def:g2:moon-in-the-sky:moon}{والقمر} الصناعي ينفق ذلك الجذب في حني مساره: فيسقط حول الأرض،
مخطئًا الهدف إلى الأبد، تمامًا كما علّم مدفع الجبل.
\textbf{8.} ``إن أُطلقت قذيفة \omterm{def:g7:motion-average-speed:formula}{بسرعة} كافية طابق هبوطها انحناء
الأرض، فدارت إلى الأبد --- ساقطةً بلا حطّ. وكل \omterm{def:g2:moon-in-the-sky:moon}{قمر} صناعي هو
تلك القذيفة، مطلَقةً جانبًا \omterm{def:g7:motion-average-speed:formula}{بسرعة} تكفي لأن تظلّ تخطئ الأرض.''
\textbf{9.} $\num{3.8e8} \div \num{6.4e6} \approx 59.4$ نصف
قطر أرضي.
\textbf{10.} نحو $1 \div 59.4^2 \approx 1/3500$ من \omterm{def:g2:pushes-and-pulls:force}{القوة}
السطحية.
\textbf{11.} رفيق لأن \omterm{def:g9:gravitation:gravitation}{الجاذبية} عند ستّين نصف قطر مخفَّفة نحو
ثلاثة آلاف وخمسمئة ضعف --- فيتقلّص هبوط الثانية الأولى البالغ
خمسة أمتار إلى نحو ملّيمتر. وبلا انتهاء لأن ذلك الملّيمتر من
السقوط يطابق تمامًا الملّيمتر الذي ينحني به سطح الأرض هابطًا
تحت انزلاق \omterm{def:g2:moon-in-the-sky:moon}{القمر} الجانبي: فالسقوط لا يكسب شيئًا أبدًا.
\textbf{12.} مثلًا: ``التفاحة \omterm{def:g2:moon-in-the-sky:moon}{والقمر} الصناعي \omterm{def:g2:moon-in-the-sky:moon}{والقمر} تطيع جملة
واحدة: كل \omterm{def:g3:measuring-mass:mass}{كتلة} تسقط نحو كل \omterm{def:g3:measuring-mass:mass}{كتلة} أخرى، بمقدار
$G m_A m_B / d^2$ --- فالتفاحة تحطّ، \omterm{def:g2:moon-in-the-sky:moon}{والقمر} الصناعي \omterm{def:g2:moon-in-the-sky:moon}{والقمر}
يظلّان يخطئان الهدف، والفرق ليس إلا \omterm{def:g7:motion-average-speed:formula}{السرعة} الجانبية.''
\end{solution}
